La verificación busca comprobar que el procesador y los componentes que lo
rodean se comporten como se especificó. Para ello se prepararon casos con
entradas y resultados esperados, que se aplicaron tanto a bloques aislados como
a conjuntos cada vez más completos. Este capítulo presenta cómo se organizaron
esas pruebas, qué se comprobó y cuáles fueron los resultados registrados.

\section{Niveles de verificación}

Una prueba puede concentrarse en un bloque particular o examinar la interacción
entre varios componentes. Los niveles permiten distinguir qué parte del
sistema se observó en cada caso:

\begin{itemize}
  \item \textbf{Modelo abstracto:} comprueba reglas de control, recuperación y
    límites mediante un modelo de su comportamiento.
  \item \textbf{Unidad:} ejercita un bloque o una lógica local de manera
    aislada, como el detector de una dependencia.
  \item \textbf{Integración local:} comprueba una etapa o la conexión entre
    varios bloques, como el decoder junto con el detector.
  \item \textbf{CPU:} ejecuta programas sobre \code{processor_top} para
    observar el funcionamiento conjunto del procesador.
  \item \textbf{Sistema:} conecta la CPU con UART, el controlador de protocolo
    y el serializador para comprobar los intercambios y sus efectos.
  \item \textbf{Cliente PC:} comprueba la aplicación y su transporte, incluidos
    el tratamiento de respuestas y los tiempos de espera.
  \item \textbf{Elaboración y revisión:} examina tipos, parámetros, interfaces
    y dependencias de la implementación.
  \item \textbf{Implementación:} utiliza Vivado y los reportes del diseño
    físico para el perfil \code{REFERENCE}.
  \item \textbf{Placa:} ejecuta pruebas en la Basys 3 y compara las respuestas
    recibidas por UART.
\end{itemize}

Estos niveles se complementan. Por ejemplo, una prueba unitaria del detector
puede comprobar que reconozca una dependencia, pero hace falta observar la CPU
para verificar cómo se retiene y avanza el pipeline. Del mismo modo, una
simulación de CPU no comprueba por sí sola la comunicación física con la placa.

Para interpretar cada resultado se registra el caso y el nivel utilizado, junto
con el estímulo, la expectativa, lo observado y las herramientas
pertinentes. Si un banco agota su tiempo de espera, se informa como problema o
limitación de esa prueba, no como una terminación \code{FINISHED} de la CPU.

\section{Manifiesto y cobertura}

Las pruebas se organizan en un manifiesto: un catálogo que identifica qué se
comprueba, con qué entradas y cuál es el resultado esperado. Una especificación
puede tener varias variantes, por ejemplo distintos operandos o direcciones
límite. Cada grupo y cada variante tienen un identificador que permite localizar
su definición y sus resultados.

El manifiesto contiene 66 grupos y 1484 casos. La
tabla~\ref{tab:manifiesto} muestra cuántos corresponden a cada área y cuántos
figuran como implementados en el ejecutor de pruebas.

\begin{table}[H]
  \centering
  \small
  \begin{tabular}{lrrrr}
    \toprule
    \textbf{Área} & \textbf{Grupos} & \textbf{Casos} & \textbf{Implementados} &
    \textbf{No implementados} \\
    \midrule
    ISA & 8 & 537 & 537 & 0 \\
    Pipeline & 17 & 249 & 249 & 0 \\
    Memoria & 10 & 200 & 199 & 1 \\
    Ejecución & 11 & 222 & 222 & 0 \\
    UART y Debug & 14 & 248 & 248 & 0 \\
    Plataforma & 6 & 28 & 25 & 3 \\
    \midrule
    Total & 66 & 1484 & 1480 & 4 \\
    \bottomrule
  \end{tabular}
  \caption{Estado de los casos del manifiesto por área.}
  \label{tab:manifiesto}
\end{table}

Un caso definido describe una situación que se quiere comprobar; una ejecución
es un intento concreto de comprobarla con un banco y un nivel determinados. El
manifiesto registra 208 configuraciones de ejecución y un mismo caso puede
utilizarse en varias de ellas, por lo que la cantidad de ejecuciones de una
campaña puede superar la cantidad de casos. Además, que un caso esté
implementado indica que existe una prueba disponible, no que haya aprobado una
corrida.

La matriz de cobertura relaciona cada requisito con su contrato y sus pruebas:
\doc{verification/coverage.md}{la matriz de cobertura}.

Los cuatro casos que el ejecutor de simulación enumera como no implementados son
\code{TEST-MEM-008/C007}, \code{TEST-PLAT-003/C001},
\code{TEST-PLAT-004/C001} y \code{TEST-PLAT-006/C001}. Sus comprobaciones se
realizan mediante otras evidencias: revisión del formato de \code{LOAD},
elaboración y revisión estructural, implementación física y revisión documental,
respectivamente. Por eso se presentan por separado de las simulaciones.

\section{Verificación funcional}

La verificación funcional comprueba tanto los resultados de las instrucciones
como el comportamiento del pipeline, las memorias y el control de sesión. Por
ejemplo, se prueban operaciones aritméticas, codificaciones del decoder,
dependencias entre instrucciones y accesos a los límites de DMEM.

Las pruebas se agrupan según el aspecto del procesador que se quiere comprobar:

\begin{itemize}
  \item \textbf{Instrucciones:} se comprueban las operaciones aritméticas y
    lógicas, la formación de inmediatos y la identificación de instrucciones
    válidas e inválidas.
  \item \textbf{Pipeline:} se observan el avance de las instrucciones, el
    forwarding, las esperas por dependencias y el descarte ante cambios de flujo.
  \item \textbf{Memorias:} se verifican los tamaños de load y store, el orden
    de los bytes, la alineación y los límites de las direcciones.
  \item \textbf{Control de ejecución:} se comprueban la preparación, los
    modos RUN y STEP, el drenado y la respuesta ante \code{halt} o faults.
\end{itemize}

Por ejemplo, el caso de \emph{load-use} presentado en el capítulo 7 comprueba
una dependencia entre un load y la instrucción siguiente. El resultado esperado
incluye tanto el valor calculado como la espera necesaria para obtenerlo. El
detalle de este y de los demás casos se encuentra en el
\doc{verification/cases.json}{manifiesto de pruebas}.

\section{Resultados de la campaña}

Una vez preparadas las pruebas, se ejecutaron simulaciones de los bloques y del
sistema completo. También se probaron programas en la placa para comprobar el
recorrido real entre la PC y la FPGA. Los resultados de la simulación y de la
placa se presentan por separado porque cada uno observa condiciones distintas.

\subsection{Regresión de simulación RTL}

La \emph{regresión} consiste en volver a ejecutar un conjunto de pruebas para
comprobar que el diseño siga cumpliendo los resultados esperados. En la corrida
general se simularon las descripciones del hardware con Icarus Verilog y
Verilator. Se realizaron 2695 ejecuciones de simulación RTL: en la primera
corrida aprobaron 2692 y las tres restantes aprobaron al repetirse.

Además de las simulaciones, se comprobaron aspectos del código y de la
infraestructura de pruebas, y se ejecutaron las 114 pruebas del assembler.
Estas comprobaciones resultaron conformes.

\subsection{Pruebas de UART y Debug}

Las pruebas de integración comprueban el intercambio entre la CPU, el
controlador de protocolo, UART y Debug. Entre los escenarios ejercitados se
encuentran la carga de programas, la transmisión de memoria utilizada y el
comportamiento del reset.

Se realizaron pruebas dirigidas de estos intercambios con la CPU y los bloques
de comunicación conectados. También se simuló la integración con los componentes
de la plataforma para observar la comunicación y el reset. El procedimiento y
los resultados se desarrollan en el capítulo~\ref{chap:integracion-sistema}.

\subsection{Validación en placa}

La validación en placa permite comprobar el recorrido completo con la FPGA y la
comunicación UART reales. Para ello se cargaron programas, se controló su
ejecución desde la PC y se compararon las respuestas con los resultados
esperados. El procedimiento, los resultados y las pruebas que quedaron fuera de
su alcance se describen en el capítulo~\ref{chap:aceptacion-fisica}.

Las instrucciones para ejecutar las pruebas se encuentran en la
\doc{verification/test_selection.md}{guía de ejecución y selección de pruebas}.
